% =============================================================
% CORIOTLAB — Plantilla Base Maestra
% Compilar con: xelatex -interaction=nonstopmode [archivo].tex
% Requiere XeLaTeX. NO usar pdflatex ni lualatex.
% =============================================================

\documentclass[12pt, a4paper, oneside]{article}

% -------------------------------------------------------------
% PAQUETES
% -------------------------------------------------------------
\usepackage{fontspec}
\usepackage{xcolor}
\usepackage{graphicx}
\usepackage{geometry}
\usepackage{fancyhdr}
\usepackage{titlesec}
\usepackage{titletoc}
\usepackage[hidelinks, unicode]{hyperref}
\usepackage{booktabs}
\usepackage{tabularx}
\usepackage{colortbl}
\usepackage{array}
\usepackage{longtable}
\usepackage{listings}
\usepackage{enumitem}
\usepackage{setspace}
\usepackage{parskip}
\usepackage{microtype}
\usepackage{eso-pic}
\usepackage{transparent}
\usepackage[most]{tcolorbox}
\usepackage{multirow}
\usepackage{caption}
\usepackage{subcaption}
\usepackage{float}
\usepackage{pdfpages}
\usepackage{lastpage}
\usepackage{calc}
\usepackage{tikz}
\usepackage{ifthen}
\usepackage{datetime2}
\usepackage[spanish, es-nodecimaldot]{babel}

% -------------------------------------------------------------
% COLORES
% -------------------------------------------------------------
% =============================================================
% CORIOTLAB — Paleta de Colores Oficial
% Brand Book v. 17 de julio de 2025
% =============================================================

\definecolor{AzulITM}     {HTML}{102D69}   % Control — C.85/M.57/Y.0/K.59
\definecolor{Magenta}     {HTML}{C14894}   % Robótica — C.0/M.63/Y.23/K.24
\definecolor{AzulDigital} {HTML}{56ACDE}   % IoT      — C.61/M.23/Y.0/K.13
\definecolor{GrisPizarra} {HTML}{2F2F2F}   % Neutro   — C.0/M.0/Y.0/K.82

% Colores auxiliares del sistema
\definecolor{GrisClaro}   {HTML}{F2F2F2}
\definecolor{GrisMedio}   {HTML}{AAAAAA}
\definecolor{GrisLinea}   {HTML}{DDDDDD}
\definecolor{Blanco}      {HTML}{FFFFFF}


% =============================================================
% ▼▼▼ ZONA DE CONFIGURACIÓN — editar antes de compilar ▼▼▼
% =============================================================

% Tipo y color del documento
\newcommand{\TipoDocumento}{informe_actividades}
\newcommand{\TipoDocumentoLabel}{INFORME DE ACTIVIDADES}
\newcommand{\ArchivoMembrete}{../../membretes/membrete_azul.png}
\newcommand{\ColorPrincipal}{AzulITM}

% Identificación del documento
\newcommand{\TituloDocumento}{Título del Informe}
\newcommand{\SubtituloDocumento}{Subtítulo o descripción breve del alcance}
\newcommand{\NumeroInforme}{001}
\newcommand{\VersionDocumento}{1.0}
\newcommand{\FechaDocumento}{\today}

% Datos del proyecto
\newcommand{\NombreProyecto}{Nombre del Proyecto}
\newcommand{\CodigoProyecto}{CL-2025-XXX}
\newcommand{\PeriodoReporte}{Mes AAAA}
\newcommand{\LugarReporte}{Medellín, Colombia}

% Autoría y destinatario
\newcommand{\AutorNombre}{Nombre Completo}
\newcommand{\AutorCargo}{Cargo del Autor}
\newcommand{\AutorEmail}{correo@coriotlab.co}
\newcommand{\DestinatarioNombre}{Nombre del Destinatario}
\newcommand{\DestinatarioCargo}{Cargo}
\newcommand{\DestinatarioEntidad}{Entidad}

% Resumen ejecutivo (portada)
\newcommand{\ResumenEjecutivo}{%
  Resumen ejecutivo del documento. Describa brevemente el
  contenido, alcance y principales hallazgos o resultados
  del informe en no más de cinco líneas.
}

% =============================================================
% ▲▲▲ FIN DE ZONA DE CONFIGURACIÓN ▲▲▲
% =============================================================

% -------------------------------------------------------------
% GEOMETRÍA
% -------------------------------------------------------------
\geometry{
  a4paper,
  top    = 3.2cm,
  bottom = 3.0cm,
  left   = 2.8cm,
  right  = 2.5cm,
  headheight = 1.2cm,
  headsep    = 0.6cm,
  footskip   = 1.0cm
}

% -------------------------------------------------------------
% FUENTES
% Rutas relativas desde la subcarpeta del documento:
%   informes/[tipo]/archivo.tex → ../../fonts/
% -------------------------------------------------------------
\setmainfont{Inter}[
  Path           = ../../fonts/Inter/static/,
  Extension      = .ttf,
  UprightFont    = *-Regular,
  BoldFont       = *-Bold,
  ItalicFont     = *-Italic,
  BoldItalicFont = *-BoldItalic
]

\newfontfamily\FuenteTitulo{MuseoModerno}[
  Path           = ../../fonts/MuseoModerno/static/,
  Extension      = .ttf,
  UprightFont    = *-Regular,
  BoldFont       = *-Bold
]

\newfontfamily\FuenteCodigo{SpaceMono}[
  Path           = ../../fonts/SpaceMono/,
  Extension      = .ttf,
  UprightFont    = *-Regular,
  BoldFont       = *-Bold,
  ItalicFont     = *-Italic,
  BoldItalicFont = *-BoldItalic
]

% -------------------------------------------------------------
% MEMBRETE DE FONDO — imagen PNG en todas las páginas
% -------------------------------------------------------------
\AddToShipoutPictureBG{%
  \begin{tikzpicture}[remember picture, overlay]
    \node[opacity=0.08, inner sep=0pt] at (current page.center) {%
      \includegraphics[width=\paperwidth, height=\paperheight,
                       keepaspectratio=false]{\ArchivoMembrete}%
    };
  \end{tikzpicture}%
}

% -------------------------------------------------------------
% ENCABEZADO Y PIE — estilo normal (páginas 2+)
% -------------------------------------------------------------
\pagestyle{fancy}
\fancyhf{}

\fancyhead[L]{%
  {\FuenteTitulo\bfseries\color{\ColorPrincipal}CORIOTLAB}%
  {\color{GrisLinea}\ |\ }%
  {\color{GrisPizarra}\TituloDocumento}%
}
\fancyhead[R]{%
  {\FuenteCodigo\small\color{GrisMedio}N°\NumeroInforme\ ·\ v\VersionDocumento}%
}
\renewcommand{\headrulewidth}{0.5pt}
\renewcommand{\headrule}{%
  {\color{\ColorPrincipal}\hrule width\headwidth height 0.5pt}%
}

\fancyfoot[L]{%
  {\FuenteCodigo\scriptsize\color{GrisMedio}%
   www.coriotlab.co\ ·\ info@coriotlab.co}%
}
\fancyfoot[C]{%
  {\scriptsize\color{GrisMedio}\FechaDocumento}%
}
\fancyfoot[R]{%
  {\FuenteTitulo\bfseries\color{\ColorPrincipal}\thepage}%
  {\color{GrisMedio}\ /\ \pageref{LastPage}}%
}
\renewcommand{\footrulewidth}{0.3pt}
\renewcommand{\footrule}{%
  {\color{GrisLinea}\hrule width\headwidth height 0.3pt}%
}

% Estilo portada — sin encabezado ni pie
\fancypagestyle{portada}{%
  \fancyhf{}%
  \renewcommand{\headrulewidth}{0pt}%
  \renewcommand{\footrulewidth}{0pt}%
}

% -------------------------------------------------------------
% ESTILOS DE SECCIÓN
% -------------------------------------------------------------
\titleformat{\section}
  {\FuenteTitulo\Large\color{\ColorPrincipal}}
  {\thesection}{0.8em}{}
  [\vspace{2pt}{\color{\ColorPrincipal}\titlerule[0.5pt]}\vspace{4pt}]

\titleformat{\subsection}
  {\FuenteTitulo\normalsize\bfseries\color{\ColorPrincipal}}
  {\thesubsection}{0.7em}{}

\titleformat{\subsubsection}
  {\FuenteTitulo\small\color{GrisPizarra}}
  {\thesubsubsection}{0.6em}{}

\titlespacing{\section}     {0pt}{18pt}{8pt}
\titlespacing{\subsection}  {0pt}{12pt}{5pt}
\titlespacing{\subsubsection}{0pt}{9pt}{3pt}

% -------------------------------------------------------------
% TABLA DE CONTENIDOS
% -------------------------------------------------------------
\contentsname{\FuenteTitulo\Large\color{\ColorPrincipal}Tabla de Contenido}

\titlecontents{section}[1.5em]
  {\vspace{4pt}\FuenteTitulo\small\color{GrisPizarra}}
  {\contentslabel{1.5em}}{\hspace*{-1.5em}}
  {\hfill\FuenteCodigo\tiny\color{GrisMedio}\contentspage}

\titlecontents{subsection}[3.5em]
  {\FuenteCodigo\footnotesize\color{GrisMedio}}
  {\contentslabel{2em}}{\hspace*{-2em}}
  {\hfill\contentspage}

% -------------------------------------------------------------
% TIPOS DE COLUMNAS DE TABLA
% -------------------------------------------------------------
\newcolumntype{C}[1]{>{\centering\arraybackslash}p{#1}}
\newcolumntype{L}[1]{>{\raggedright\arraybackslash}p{#1}}
\newcolumntype{R}[1]{>{\raggedleft\arraybackslash}p{#1}}

% Macro encabezado de tabla con color principal
\newcommand{\encabezadoTabla}{\rowcolor{\ColorPrincipal}\color{white}\FuenteTitulo\small\bfseries}

% -------------------------------------------------------------
% ESTILOS DE CÓDIGO
% -------------------------------------------------------------
\lstset{
  basicstyle   = \FuenteCodigo\small,
  backgroundcolor = \color{GrisClaro},
  keywordstyle = \bfseries\color{\ColorPrincipal},
  commentstyle = \itshape\color{GrisMedio},
  stringstyle  = \color{Magenta},
  numberstyle  = \FuenteCodigo\tiny\color{GrisMedio},
  numbers      = left,
  numbersep    = 8pt,
  frame        = single,
  rulecolor    = \color{GrisLinea},
  breaklines   = true,
  showstringspaces = false,
  tabsize      = 2,
  captionpos   = b,
  xleftmargin  = 1.5em,
  xrightmargin = 0.5em,
  aboveskip    = 10pt,
  belowskip    = 6pt
}

\lstdefinelanguage{Python3}{
  language     = Python,
  morekeywords = {True, False, None, async, await, f}
}

% -------------------------------------------------------------
% CAJAS TCOLORBOX
% -------------------------------------------------------------
\tcbuselibrary{skins, breakable}

\newtcolorbox{cajaResultado}[1][]{
  colframe    = \ColorPrincipal,
  colback     = GrisClaro,
  colbacktitle= \ColorPrincipal,
  coltitle    = white,
  title       = {{\FuenteTitulo\small\bfseries Resultado}},
  arc         = 3pt,
  breakable,
  enhanced,
  #1
}

\newtcolorbox{cajaNota}[1][]{
  colframe    = AzulDigital,
  colback     = GrisClaro,
  colbacktitle= AzulDigital,
  coltitle    = white,
  title       = {{\FuenteTitulo\small\bfseries Nota}},
  arc         = 3pt,
  breakable,
  enhanced,
  #1
}

\newtcolorbox{cajaAlerta}[1][]{
  colframe    = Magenta,
  colback     = GrisClaro,
  colbacktitle= Magenta,
  coltitle    = white,
  title       = {{\FuenteTitulo\small\bfseries Importante}},
  arc         = 3pt,
  breakable,
  enhanced,
  #1
}

\newtcolorbox{cajaDatos}[1][]{
  colframe    = GrisPizarra,
  colback     = GrisClaro,
  colbacktitle= GrisPizarra,
  coltitle    = white,
  title       = {{\FuenteTitulo\small\bfseries Información}},
  arc         = 3pt,
  breakable,
  enhanced,
  #1
}

% -------------------------------------------------------------
% ESPACIADO
% -------------------------------------------------------------
\setlength{\parskip}{6pt}
\setlength{\parindent}{0pt}
\onehalfspacing

% =============================================================
% DOCUMENTO
% =============================================================
\begin{document}

% -------------------------------------------------------------
% PORTADA
% -------------------------------------------------------------
\thispagestyle{portada}

\vspace*{2.5cm}

% Franja de tipo de documento
\begin{tikzpicture}
  \fill[\ColorPrincipal] (0,0) rectangle (\textwidth, 1cm);
  \node[anchor=west, xshift=10pt, yshift=0.5cm] at (0,0) {%
    {\FuenteTitulo\small\bfseries\color{white}\TipoDocumentoLabel}%
  };
  \node[anchor=east, xshift=-10pt, yshift=0.5cm] at (\textwidth,0) {%
    {\FuenteCodigo\small\color{white}N°\NumeroInforme}%
  };
\end{tikzpicture}

\vspace{1.2cm}

{\FuenteTitulo\huge\color{\ColorPrincipal}\TituloDocumento\par}
\vspace{0.4cm}
{\FuenteTitulo\large\color{GrisPizarra}\SubtituloDocumento\par}
\vspace{0.6cm}

{\color{AzulDigital}\hrule height 0.4pt}
\vspace{0.8cm}

% Tabla de metadatos
\begin{tabular}{L{3.2cm} L{10cm}}
  {\FuenteCodigo\footnotesize\color{GrisMedio}PROYECTO}       & {\small\NombreProyecto} \\[4pt]
  {\FuenteCodigo\footnotesize\color{GrisMedio}CÓDIGO}         & {\small\CodigoProyecto} \\[4pt]
  {\FuenteCodigo\footnotesize\color{GrisMedio}PERÍODO}        & {\small\PeriodoReporte} \\[4pt]
  {\FuenteCodigo\footnotesize\color{GrisMedio}LUGAR}          & {\small\LugarReporte} \\[4pt]
  {\FuenteCodigo\footnotesize\color{GrisMedio}ELABORADO POR}  & {\small\AutorNombre\ —\ \AutorCargo} \\[4pt]
  {\FuenteCodigo\footnotesize\color{GrisMedio}DIRIGIDO A}     & {\small\DestinatarioNombre\ —\ \DestinatarioCargo,\ \DestinatarioEntidad} \\[4pt]
  {\FuenteCodigo\footnotesize\color{GrisMedio}VERSIÓN}        & {\small v\VersionDocumento} \\[4pt]
  {\FuenteCodigo\footnotesize\color{GrisMedio}FECHA}          & {\small\FechaDocumento} \\
\end{tabular}

\vspace{0.8cm}

\begin{tcolorbox}[
  colframe    = \ColorPrincipal,
  colback     = GrisClaro,
  colbacktitle= \ColorPrincipal,
  coltitle    = white,
  title       = {{\FuenteTitulo\small\bfseries Resumen Ejecutivo}},
  arc         = 3pt,
  enhanced
]
\ResumenEjecutivo
\end{tcolorbox}

\newpage

% -------------------------------------------------------------
% TABLA DE CONTENIDO
% -------------------------------------------------------------
\tableofcontents
\newpage

% =============================================================
% CUERPO DEL DOCUMENTO
% Descomenta los bloques que necesitas según el tipo de informe
% =============================================================

% ─────────────────────────────────────────────────────────────
% INTRODUCCIÓN (presente en todos los documentos)
% ─────────────────────────────────────────────────────────────
\section{Introducción}

Describa aquí el contexto general del documento, su propósito
y el alcance de las actividades o contenidos presentados.
Indique brevemente el período cubierto y los sistemas o
proyectos involucrados.

% ─────────────────────────────────────────────────────────────
% [BLOQUE A — Objetivos]
% Uso: actividades, tecnico, propuesta
% ─────────────────────────────────────────────────────────────
\section{Objetivos}

\subsection{Objetivo General}
Describa el objetivo principal que orienta el trabajo
documentado en este informe.

\subsection{Objetivos Específicos}
\begin{itemize}[label={\color{\ColorPrincipal}>}, leftmargin=1.5em]
  \item Objetivo específico 1.
  \item Objetivo específico 2.
  \item Objetivo específico 3.
\end{itemize}

% ─────────────────────────────────────────────────────────────
% [BLOQUE B — Actividades Realizadas]
% Uso: actividades, tecnico
% ─────────────────────────────────────────────────────────────
\section{Actividades Realizadas}

\subsection{Descripción de Actividades}

A continuación se presenta el resumen de las actividades
ejecutadas durante el período reportado.

\vspace{6pt}
\begin{longtable}{C{0.5cm} L{3.8cm} L{5.5cm} C{1.2cm} C{1.8cm}}
  \toprule
  \encabezadoTabla N° &
  \encabezadoTabla Actividad &
  \encabezadoTabla Descripción &
  \encabezadoTabla Horas &
  \encabezadoTabla Estado \\
  \midrule
  \endfirsthead

  \multicolumn{5}{l}{\footnotesize\itshape (Continúa de la página anterior)}\\
  \toprule
  \encabezadoTabla N° &
  \encabezadoTabla Actividad &
  \encabezadoTabla Descripción &
  \encabezadoTabla Horas &
  \encabezadoTabla Estado \\
  \midrule
  \endhead

  \midrule
  \multicolumn{5}{r}{\footnotesize\itshape (Continúa en la siguiente página)}\\
  \endfoot

  \bottomrule
  \endlastfoot

  \rowcolor{GrisClaro}
  1 & Nombre actividad 1 & Descripción detallada & 4 & Completado \\
  2 & Nombre actividad 2 & Descripción detallada & 6 & Completado \\
  \rowcolor{GrisClaro}
  3 & Nombre actividad 3 & Descripción detallada & 3 & En proceso \\
  \midrule
  \multicolumn{3}{r}{\FuenteTitulo\small\bfseries\color{\ColorPrincipal}Total} &
  \multicolumn{2}{l}{\bfseries 13 horas} \\
\end{longtable}

% ─────────────────────────────────────────────────────────────
% [BLOQUE C — Marco Teórico]
% Uso: tecnico, propuesta
% Comentar si no aplica
% ─────────────────────────────────────────────────────────────
%\section{Marco Teórico}
%
%\subsection{Fundamento Conceptual}
%Describa los conceptos teóricos que sustentan el trabajo
%desarrollado, incluyendo principios de control, robótica o IoT
%según corresponda al área del proyecto.
%
%\subsection{Tecnologías Utilizadas}
%Enumere y describa brevemente las tecnologías, frameworks,
%protocolos y plataformas empleadas en el desarrollo.

% ─────────────────────────────────────────────────────────────
% [BLOQUE D — Desarrollo Técnico + Código]
% Uso: tecnico
% Comentar si no aplica
% ─────────────────────────────────────────────────────────────
%\section{Desarrollo Técnico}
%
%\subsection{Arquitectura del Sistema}
%Describa la arquitectura general de la solución, incluyendo
%diagrama de bloques si aplica.
%
%\subsection{Implementación}
%Detalle los componentes implementados y su integración.
%
% Ejemplo: lectura de sensor IoT en Python
%\begin{lstlisting}[language=Python3, caption={Lectura de sensor IoT}]
%import machine
%import time
%
%# Configuracion del sensor
%sensor_pin = machine.ADC(machine.Pin(34))
%sensor_pin.atten(machine.ADC.ATTN_11DB)
%
%def leer_sensor():
%    """Lee el valor analogico del sensor y convierte a voltaje."""
%    valor_raw = sensor_pin.read()
%    voltaje = (valor_raw / 4095.0) * 3.3
%    return voltaje
%
%while True:
%    v = leer_sensor()
%    print(f"Voltaje: {v:.2f} V")
%    time.sleep(1)
%\end{lstlisting}
%
% Ejemplo: publicacion MQTT en C++ para ESP32
%\begin{lstlisting}[language=C++, caption={Publicacion MQTT — ESP32}]
%#include <WiFi.h>
%#include <PubSubClient.h>
%
%const char* ssid     = "RED_WIFI";
%const char* password = "CONTRASENA";
%const char* mqtt_server = "192.168.1.100";
%
%WiFiClient espClient;
%PubSubClient client(espClient);
%
%void setup() {
%  Serial.begin(115200);
%  WiFi.begin(ssid, password);
%  while (WiFi.status() != WL_CONNECTED) delay(500);
%  client.setServer(mqtt_server, 1883);
%}
%
%void loop() {
%  if (!client.connected()) {
%    client.connect("ESP32_CORIOT");
%  }
%  // Publica temperatura simulada
%  client.publish("coriot/sensor/temp", "25.3");
%  delay(2000);
%}
%\end{lstlisting}

% ─────────────────────────────────────────────────────────────
% [BLOQUE E — Metodología y Cronograma]
% Uso: propuesta
% Comentar si no aplica
% ─────────────────────────────────────────────────────────────
%\section{Metodología}
%
%Describa el enfoque metodológico empleado o propuesto para
%el desarrollo del proyecto.
%
%\subsection{Cronograma}
%
%\begin{longtable}{C{0.5cm} L{2.8cm} L{5.0cm} C{1.8cm} C{1.8cm}}
%  \toprule
%  \encabezadoTabla N° &
%  \encabezadoTabla Fase &
%  \encabezadoTabla Descripción &
%  \encabezadoTabla Inicio &
%  \encabezadoTabla Fin \\
%  \midrule
%  \endfirsthead
%  \toprule
%  \encabezadoTabla N° &
%  \encabezadoTabla Fase &
%  \encabezadoTabla Descripción &
%  \encabezadoTabla Inicio &
%  \encabezadoTabla Fin \\
%  \midrule
%  \endhead
%  \bottomrule
%  \endlastfoot
%  \rowcolor{GrisClaro}
%  1 & Diagnóstico & Levantamiento de requerimientos & 01/07/2025 & 15/07/2025 \\
%  2 & Diseño & Arquitectura y especificaciones & 16/07/2025 & 31/07/2025 \\
%  \rowcolor{GrisClaro}
%  3 & Desarrollo & Implementación del sistema & 01/08/2025 & 30/09/2025 \\
%  4 & Pruebas & Validación y ajustes & 01/10/2025 & 15/10/2025 \\
%  \rowcolor{GrisClaro}
%  5 & Entrega & Documentación y cierre & 16/10/2025 & 31/10/2025 \\
%\end{longtable}

% ─────────────────────────────────────────────────────────────
% [BLOQUE F — Presupuesto y Pago]
% Uso: propuesta, cuenta_cobro
% Comentar si no aplica
% ─────────────────────────────────────────────────────────────
%\section{Presupuesto}
%
%\begin{longtable}{C{0.5cm} L{3.0cm} L{4.5cm} C{1.5cm} R{2.5cm}}
%  \toprule
%  \encabezadoTabla N° &
%  \encabezadoTabla Ítem &
%  \encabezadoTabla Descripción &
%  \encabezadoTabla Cantidad &
%  \encabezadoTabla Valor (COP) \\
%  \midrule
%  \endfirsthead
%  \toprule
%  \encabezadoTabla N° &
%  \encabezadoTabla Ítem &
%  \encabezadoTabla Descripción &
%  \encabezadoTabla Cantidad &
%  \encabezadoTabla Valor (COP) \\
%  \midrule
%  \endhead
%  \bottomrule
%  \endlastfoot
%  \rowcolor{GrisClaro}
%  1 & Honorarios & Asesoría técnica IoT & 1 & \$1.500.000 \\
%  2 & Materiales & Componentes electrónicos & 5 & \$350.000 \\
%  \rowcolor{GrisClaro}
%  3 & Software & Licencias y herramientas & 1 & \$200.000 \\
%  \midrule
%  \multicolumn{4}{r}{\small Subtotal} & \$2.050.000 \\
%  \multicolumn{4}{r}{\small IVA (19\%)} & \$389.500 \\
%  \rowcolor{GrisClaro}
%  \multicolumn{4}{r}{\bfseries\FuenteTitulo\color{\ColorPrincipal}TOTAL} &
%    {\bfseries \$2.439.500} \\
%\end{longtable}
%
%\begin{cajaDatos}[title={{\FuenteTitulo\small\bfseries Datos Bancarios para Pago}}]
%\begin{tabular}{L{3cm} L{8cm}}
%  \textbf{Banco:}          & Banco de Bogotá \\
%  \textbf{Tipo de cuenta:} & Ahorros \\
%  \textbf{N° cuenta:}      & 000-000000-0 \\
%  \textbf{Titular:}        & \AutorNombre \\
%  \textbf{NIT / CC:}       & 0.000.000.000 \\
%  \textbf{Contacto:}       & \AutorEmail \\
%\end{tabular}
%\end{cajaDatos}

% ─────────────────────────────────────────────────────────────
% [BLOQUE G — Resultados]
% Uso: actividades, tecnico, propuesta
% ─────────────────────────────────────────────────────────────
\section{Resultados}

\begin{cajaResultado}
Describa aquí los resultados obtenidos durante el período,
destacando los logros más relevantes frente a los objetivos
planteados al inicio del informe.
\end{cajaResultado}

\vspace{8pt}

\subsection{Indicadores de Cumplimiento}

\begin{tabular}{L{3.5cm} L{4.0cm} C{1.8cm} C{1.8cm}}
  \toprule
  \encabezadoTabla Indicador &
  \encabezadoTabla Descripción &
  \encabezadoTabla Meta &
  \encabezadoTabla Logro \\
  \midrule
  \rowcolor{GrisClaro}
  Indicador 1 & Descripción del indicador & 100\% & 95\% \\
  Indicador 2 & Descripción del indicador & 5 & 5 \\
  \rowcolor{GrisClaro}
  Indicador 3 & Descripción del indicador & 3 & 2 \\
  \bottomrule
\end{tabular}

\vspace{8pt}

% Espacio para figura — reemplazar \fbox con \includegraphics
% \begin{figure}[H]
%   \centering
%   \fbox{\parbox{10cm}{\centering\vspace{3cm}Insertar figura aquí\vspace{3cm}}}
%   \caption{Descripción de la figura}
%   \label{fig:resultado}
% \end{figure}

% ─────────────────────────────────────────────────────────────
% [BLOQUE H — Conclusiones]
% Uso: actividades, tecnico, propuesta, memorando
% ─────────────────────────────────────────────────────────────
\section{Conclusiones}

\begin{itemize}[label={\color{\ColorPrincipal}>}, leftmargin=1.5em]
  \item Conclusión principal del período reportado.
  \item Aspectos positivos identificados durante la ejecución.
  \item Dificultades encontradas y cómo se resolvieron.
  \item Evaluación general del cumplimiento de objetivos.
\end{itemize}

% ─────────────────────────────────────────────────────────────
% [BLOQUE I — Recomendaciones]
% Uso: tecnico, propuesta
% Comentar si no aplica
% ─────────────────────────────────────────────────────────────
%\section{Recomendaciones}
%
%\begin{cajaNota}[title={{\FuenteTitulo\small\bfseries Próximos Pasos}}]
%\begin{itemize}[label={\color{AzulDigital}>}, leftmargin=1.5em]
%  \item Recomendación 1 para el siguiente período.
%  \item Recomendación 2 sobre aspectos técnicos o metodológicos.
%  \item Recomendación 3 sobre recursos o colaboraciones.
%\end{itemize}
%\end{cajaNota}

% ─────────────────────────────────────────────────────────────
% [BLOQUE J — Referencias]
% Uso: tecnico, propuesta
% Comentar si no aplica
% ─────────────────────────────────────────────────────────────
%\section{Referencias}
%
%\begin{enumerate}[label={[\arabic*]}, leftmargin=2em]
%  \item Apellido, N. (Año). \textit{Título del documento}.
%        Editorial / Fuente. URL si aplica.
%  \item Apellido, N. (Año). \textit{Título del artículo}.
%        \textit{Revista}, Vol(N°), pp. 00-00.
%  \item IEEE Std 802.11-2020. \textit{IEEE Standard for
%        Information Technology}. IEEE.
%\end{enumerate}

% ─────────────────────────────────────────────────────────────
% [BLOQUE FIJO — Firma y Aprobación]
% Presente en TODOS los documentos — NO comentar
% ─────────────────────────────────────────────────────────────
\newpage
\section*{Firma y Aprobación}
\addcontentsline{toc}{section}{Firma y Aprobación}

\vspace{0.6cm}

\begin{tabular}{p{0.46\textwidth} p{0.08\textwidth} p{0.46\textwidth}}
  \begin{minipage}[t]{0.46\textwidth}
    \vspace{2.2cm}
    {\color{\ColorPrincipal}\hrule height 0.5pt}
    \vspace{4pt}
    {\bfseries\AutorNombre}\\
    {\small\color{GrisMedio}\AutorCargo}\\
    {\FuenteCodigo\footnotesize\color{GrisMedio}\AutorEmail}\\
    \vspace{4pt}
    {\FuenteCodigo\footnotesize\color{GrisMedio}Elaborado por}
  \end{minipage}
  & &
  \begin{minipage}[t]{0.46\textwidth}
    \vspace{2.2cm}
    {\color{\ColorPrincipal}\hrule height 0.5pt}
    \vspace{4pt}
    {\bfseries\DestinatarioNombre}\\
    {\small\color{GrisMedio}\DestinatarioCargo}\\
    {\small\color{GrisMedio}\DestinatarioEntidad}\\
    \vspace{4pt}
    {\FuenteCodigo\footnotesize\color{GrisMedio}Revisado y aprobado por}
  \end{minipage}
\end{tabular}

% ─────────────────────────────────────────────────────────────
% [BLOQUE OPCIONAL — Anexos]
% Uso: tecnico (comentado por defecto)
% ─────────────────────────────────────────────────────────────
%\newpage
%\appendix
%\section{Anexos}
%
% Para incluir un PDF externo como anexo:
%\includepdf[pages=-, pagecommand={}]{../../assets/imagenes/anexo.pdf}

\end{document}
